%!TEX program = xelatex
% Portfolio Optimisation with Causal Networks — 18-minute MSc presentation
% Built on the Imperial College London beamer template (compile with XeLaTeX).
% Structure mirrors the thesis: Introduction; Background and related work; Methodology (three chapters);
% Implementation; Results; Conclusions.
% Figures expected in ./figures/ : phase_ii_window_gradient.png

\documentclass[
	aspectratio=169,
	t,
	onlytextwidth,
	10pt,
]{beamer}

\usetheme{Imperial}

\usepackage{tikz}
\usetikzlibrary{positioning,calc,arrows.meta,shapes.misc,matrix,fit}
\usepackage{booktabs}
\usepackage{amsmath}
\usepackage{multirow}
\usepackage{adjustbox}

%----------------------------------------------------------------------------------------
%	LOCAL MACROS AND COLOURS
%----------------------------------------------------------------------------------------

% Palette matched to the report's Figure 5.1
\definecolor{skel}{HTML}{E0A317}    % skeleton effect (gold)
\definecolor{orient}{HTML}{5A3593}  % orientation effect (purple)
\definecolor{ctrl}{HTML}{6E6E6E}    % controls / neutral grey

% Allocator names, always the report's naming
\newcommand{\al}[1]{{\textsf{#1}}}
\newcommand{\corr}{\al{CORRELATION-HRP}}
\newcommand{\skl}{\al{SKELETON-HRP}}
\newcommand{\und}{\al{UNDIRECTED-HRP}}
\newcommand{\sem}{\al{SEMCOV-HRP}}
\newcommand{\topo}{\al{TOPO-HRP}}
\newcommand{\toposem}{\al{TOPO-SEMCOV-HRP}}
\newcommand{\lw}{\al{LW-HRP}}
\newcommand{\dfc}{\al{DEFACTORED-HRP}}
\newcommand{\pc}{\al{PCORR-HRP}}

\newcommand{\mn}[1]{$-#1$}
\newcommand{\pos}[1]{$+#1$}
\newcommand{\pv}[1]{{\scriptsize\textcolor{ctrl}{(#1)}}}
\newcommand{\key}[1]{\textcolor{ICLBlue}{\textbf{#1}}}
\newcommand{\pval}[1]{{\tiny\textcolor{ctrl}{(#1)}}}

% Dense slides: tighter paragraph spacing than the theme's 10pt default
% The theme re-applies its 10pt parskip inside every column (via \@minipagerestore), so override both.
\makeatletter
\newcommand{\dense}{\small\setlength{\parskip}{4pt}\def\@minipagerestore{\setlength{\parskip}{4pt}}}
\makeatother

\setlength{\leftmargini}{1.2em}

% References, set in beamer's footnote area (which every theme reserves above its footline).
% \cn{n} is a superscript citation marker on the slide; the same \cn{n} labels the entry in the footer.
% Numbering is cumulative across the deck in reading order (down the left column, then the right): a work gets a
% number the first time it is cited and keeps it thereafter. Each footer lists only the works cited for the first time on
% that slide; a repeated \cn{n} in the body refers back to the footer where n first appeared. The numbers do not match the thesis bibliography.
\newcommand{\cn}[1]{\textsuperscript{\textcolor{ICLBlue}{#1}}}
\setbeamertemplate{footnote}{\insertfootnotetext}
\newcommand{\refs}[1]{\ifx\\#1\\\else\footnotetext{\tiny\textcolor{ctrl}{#1}}\fi}

% Footline: the theme's layout, with the slide number bottom right in place of the date.
\makeatletter
\setbeamertemplate{footline}{
	\begin{beamercolorbox}[ht=2mm, dp=5mm]{page number in head/foot}
		\usebeamerfont{page number in head/foot}%
		\hspace{\beamer@leftmargin}%
		\textbf{Imperial College London}\hspace{0.5cm}\inserttitle\hfill\insertframenumber%
		\hspace*{\beamer@rightmargin}% starred: TeX discards ordinary glue at the end of a line, which pushed the number into the edge
	\end{beamercolorbox}
}
\makeatother

%----------------------------------------------------------------------------------------
%	PRESENTATION INFORMATION
%----------------------------------------------------------------------------------------

\title{Portfolio Optimisation with Causal Networks}
% \subtitle{\normalsize When the correlation matrix in hierarchical portfolio allocation is replaced by a
% causal graph, is there a performance improvement, and how much of any improvement is attributable to the graph's skeleton versus its edge orientations?}
\author{Joshua Muthu — 06048741}
\date{September 2026}

\begin{document}

%----------------------------------------------------------------------------------------
%	TITLE
%----------------------------------------------------------------------------------------

\begingroup
	\setbeamercolor{background canvas}{bg=ICLBlue}
	\setbeamercolor{title page title}{fg=white}
	\setbeamercolor{title page subtitle}{fg=white}
	\setbeamercolor{author}{fg=white}
	\setbeamercolor{date}{fg=white}
	\setbeamertemplate{title page}[logo]{ICL_Logo_White.pdf}
	\frame[plain, s]{\titlepage}
\endgroup

%----------------------------------------------------------------------------------------
%	1. INTRODUCTION — MOTIVATION
%----------------------------------------------------------------------------------------

\begin{frame}
	\frametitle{Motivation: correlation is undirected; HRP requires symmetric inputs}
	\dense
	\begin{columns}[T]
		\begin{column}{0.52\linewidth}
			\textbf{Hierarchical Risk Parity}\cn{1}\\[1em]
            \begin{tikzpicture}[
            	box/.style={draw=ICLBlue, rounded corners=2pt, text width=2.2cm, align=center, minimum height=1.1cm, font=\scriptsize, fill=ICLLightGrey!60},
            	arr/.style={-{Stealth[length=5pt]}, ICLBlue, thick}]
            	\node[box] (a) {\textbf{1. Cluster}\\ on $\sqrt{(1-\rho_{ij})/2}$};
            	\node[box, right=0.25cm of a] (b) {\textbf{2. Order}\\ by leaves};
            	\node[box, right=0.25cm of b] (c) {\textbf{3. Bisect}\\ inverse-variance splits};
            	\draw[arr] (a) -- (b); \draw[arr] (b) -- (c);
            	\node[below=0.1cm of a, font=\tiny, text=ctrl] {needs a symmetric distance};
            	\node[below=0.1cm of c, font=\tiny, text=ctrl] {needs ordering and covariance};
            \end{tikzpicture}\\[1em]
			\begin{itemize}
				\setlength{\itemsep}{3pt}
				\item Mean-variance optimisation\cn{2} is ill-conditioned and unstable.
                \item HRP avoids this instability by never inverting the covariance matrix.
				\item But correlation can be spurious\cn{3}, and records co-movement, not causal influence.
			\end{itemize}
		\end{column}
		\begin{column}{0.44\linewidth}
			\textbf{Causal discovery}
			\begin{itemize}
				\setlength{\itemsep}{3pt}
				\item Causal discovery algorithms (NOTEARS\cn{4}, DYNOTEARS\cn{5}) fit DAGs, giving directed causal links between assets.
				\item Applied to equities before\cn{6}, but not within hierarchical allocators.
				\item HRP's clustering requires a symmetric matrix: DAGs' \textcolor{skel}{\textbf{skeletons}} are symmetric, but not their \textcolor{orient}{\textbf{edge orientations}}.
				\item So the two components have to be input differently and their values measured separately.
			\end{itemize}
		\end{column}
	\end{columns}
	\refs{\cn{1} López de Prado (2016); \cn{2} Markowitz (1952); \cn{3} Simon (1954); \cn{4} Zheng et al.\ (2018); \cn{5} Pamfil et al.\ (2020); \cn{6} Howard, Lohre \& Mudde (2025)}
\end{frame}

%----------------------------------------------------------------------------------------
%	1. INTRODUCTION — RESEARCH QUESTION
%----------------------------------------------------------------------------------------

\begin{frame}
	\frametitle{Research question}
	\dense
    \vspace{-1em}
    \setbeamercolor{block body}{bg=ICLBlue, fg=white}
	\begin{block}{}
		\ImperialSansSemiBold When the correlation matrix in hierarchical portfolio allocation is replaced by a causal graph, is there a performance improvement, and how much of any improvement is attributable to the graph's skeleton versus its edge orientations?
	\end{block}
	\begin{columns}[T]
		\begin{column}{0.40\linewidth}
			\vspace{4pt}
            \hspace*{5mm}
			\begin{tikzpicture}[scale=0.85, transform shape,
				n/.style={circle, draw=ICLBlue, fill=white, inner sep=1pt, font=\tiny, minimum size=0.4cm},
				e/.style={-{Stealth[length=4pt]}, ICLBlue, thick},
				lab/.style={font=\scriptsize\bfseries, align=center}]
				\node[n] (a) at (0,1.1) {A}; \node[n] (b) at (0.9,1.6) {B}; \node[n] (c) at (0.9,0.6) {C};
				\node[n] (d) at (1.8,1.1) {D}; \node[n] (f) at (1.8,0.1) {E};
				\draw[e] (a)--(b); \draw[e] (a)--(c); \draw[e] (b)--(d); \draw[e] (c)--(d); \draw[e] (c)--(f);
				\node[lab, text=ICLBlue] at (0.9,2.1) {DAG};
				\draw[-{Stealth[length=5pt]}, ctrl, thick] (2.4,1.5) -- (3.2,2.1);
				\draw[-{Stealth[length=5pt]}, ctrl, thick] (2.4,0.5) -- (3.2,-0.2);
				\begin{scope}[shift={(3.6,1.6)}, every node/.style={n, draw=skel}]
					\node (a2) at (0,0.6) {A}; \node (b2) at (0.7,1.0) {B}; \node (c2) at (0.7,0.2) {C}; \node (d2) at (1.4,0.6) {D}; \node (f2) at (1.4,-0.3) {E};
					\draw[skel, thick] (a2)--(b2) (a2)--(c2) (b2)--(d2) (c2)--(d2) (c2)--(f2);
				\end{scope}
				\node[lab, text=skel] at (4.3,3.2) {skeleton\\ \tiny undirected structure};
				\begin{scope}[shift={(3.6,-1.2)}]
					\draw[-{Stealth[length=4pt]}, orient, thick] (0.15,0.7)--(0.6,0.95);
					\draw[-{Stealth[length=4pt]}, orient, thick] (0.15,0.5)--(0.6,0.3);
					\draw[-{Stealth[length=4pt]}, orient, thick] (0.85,0.95)--(1.3,0.7);
					\draw[-{Stealth[length=4pt]}, orient, thick] (0.85,0.3)--(1.3,0.55);
					\draw[-{Stealth[length=4pt]}, orient, thick] (0.85,0.15)--(1.3,-0.15);
				\end{scope}
				\node[lab, text=orient] at (4.4,0.2) {orientations\\ \tiny directions of edges};
			\end{tikzpicture}
		\end{column}
		\begin{column}{0.56\linewidth}
			\begin{itemize}
				\setlength{\itemsep}{3pt}
				\item Sub-questions on the mechanisms: 
                \begin{itemize}
                    \item Is the \textcolor{skel}{\textbf{skeleton's}} effect specific to the discovered graph, or is it replicated by
                    any sparse network of similar density?
                    \item Is the \textcolor{orient}{\textbf{edge orientations'}} effect due to directionality, or to some other property of the implied covariance?
                \end{itemize}
				\item Three challenges:
                \begin{enumerate}
                  \item Attributing performance to the graph's components
                        (needs a credible, reproducible graph).
                  \item Getting an asymmetric matrix into HRP.
                  \item Inference at realistic effect sizes.
                \end{enumerate}
			\end{itemize}
		\end{column}
	\end{columns}
    \vspace{0.1em}
    \key{Answer:} no performance improvement. The \textcolor{skel}{\textbf{skeleton}} has a negative effect, \textcolor{orient}{\textbf{orientations}} have a positive effect, but neither effect needs the graph.
	\refs{}
\end{frame}

%----------------------------------------------------------------------------------------
%	2. BACKGROUND AND RELATED WORK
%----------------------------------------------------------------------------------------

\begin{frame}
	\frametitle{Prior research: no direct allocation from a discovered causal graph}
	\dense
	\begin{columns}[T]
		\begin{column}{0.50\linewidth}
			\textbf{Hierarchical allocation and backtest overfitting}
			\begin{itemize}
				\setlength{\itemsep}{2pt}
				\item HRP\cn{1}, HERC\cn{7}, HSP\cn{8}.
				\item Backtest overfitting and multiple comparisons problem: Deflated Sharpe\cn{9}, Hansen's SPA\cn{10}, Model Confidence Set\cn{11}.
			\end{itemize}
			\textbf{Causal graphs in investing}
			\begin{itemize}
				\setlength{\itemsep}{2pt}
				\item Howard et al.\cn{6} apply DYNOTEARS to the S\&P 500. The contemporaneous matrix is used for similarity and centrality, but not portfolio weights.
			\end{itemize}
		\end{column}
		\begin{column}{0.46\linewidth}
			\textbf{The gap}\\[4pt]
			{\scriptsize\adjustbox{max width=\linewidth}{
			\begin{tabular}{@{}l c c c c@{}}
				\toprule
				 & discovery & weights & asymm. & decomp. \\
				\midrule
				HRP, HSP & -- & -- & -- & -- \\
				MST / PMFG & -- & \checkmark & -- & -- \\
				Granger nets & -- & \checkmark & \checkmark & -- \\
				Howard et al.\cn{6} & \checkmark & -- & -- & -- \\
				\textbf{This thesis} & \checkmark & \checkmark & \checkmark & \checkmark \\
				\bottomrule
			\end{tabular}}}\\[4pt]
			\textbf{Causal discovery methods}\\[4pt]
			\begin{itemize}
				\item NOTEARS/DYNOTEARS\cn{5} fit a linear SEM with an acyclicality constraint to variables, which scales well.
				\item VARLiNGAM\cn{12} identifies DAGs by assessing the independence of non-Gaussian residuals.
			\end{itemize}
		\end{column}
	\end{columns}
	\refs{\cn{7} Raffinot (2018); \cn{8} Rodríguez Domínguez (2023); \cn{9} Bailey \& López de Prado (2014); \cn{10} Hansen (2005); \cn{11} Hansen et al.\ (2011); \cn{12} Hyvärinen et al.\ (2010)}
\end{frame}

%----------------------------------------------------------------------------------------
%	3. METHODOLOGY I — DISCOVERING THE GRAPH
%----------------------------------------------------------------------------------------

\begin{frame}
	\frametitle{Discovering the causal graph over the assets and drivers}
	\dense
	\begin{columns}[T]
		\begin{column}{0.50\linewidth}
			\textbf{Data}
			\begin{itemize}
				\setlength{\itemsep}{3pt}
				\item 99 US large-cap assets (S\&P 100 proxy), CRSP via WRDS, fixed 2007--2024. Survivorship bias (excluding dropouts) affects levels, not pairwise comparisons.
				\item 33 plausibly exogenous macro drivers (rates, credit, FX, commodities, VIX), standardised per window.
				\item 4 estimation windows: 189, 252, 378, and 504 days; 215 monthly rebalances.
			\end{itemize}
			\textbf{Causal discovery algorithms}
			\begin{itemize}
				\setlength{\itemsep}{3pt}
				\item Primarily \key{DYNOTEARS}\cn{5} ($p=1$, $\lambda=0.05$).
				\item \key{VARLiNGAM}\cn{12,13} comparator; \key{ridge-Granger}\cn{14} as ancillary control on edge magnitudes.
			\end{itemize}
		\end{column}
		\begin{column}{0.46\linewidth}
			\begin{center}
            \vspace{-1em}
			\begin{tikzpicture}[
				cell/.style={draw=ICLBlue, minimum width=1.85cm, minimum height=1.0cm, font=\scriptsize, align=center, inner xsep=1pt},
				lab/.style={font=\scriptsize, text=ctrl}]
				\node[cell, fill=ICLLightGrey!60] (dd) at (0,1.0) {drivers$\to$drivers};
				\node[cell, fill=ICLLightGrey!60] (da) at (1.85,1.0) {drivers$\to$assets};
				\node[cell, fill=white, text=ctrl] (ad) at (0,0) {assets$\to$drivers\\ \textbf{forbidden}};
				\node[cell, fill=ICLCream, very thick] (aa) at (1.85,0) {assets$\to$assets\\ $= M$};
				\node[lab, above=0.05cm of dd] {$D$}; \node[lab, above=0.05cm of da] {$A$};
				\node[lab, left=0.05cm of dd] {$D$}; \node[lab, left=0.05cm of ad] {$A$};
			\end{tikzpicture}
			\end{center}
			\begin{itemize}
				\setlength{\itemsep}{3pt}
				\item $M$: contemporaneous asset$\to$asset block of the DAG, net of drivers' effects. Edge density: 6--7\% for DYNOTEARS, 32\% for VARLiNGAM at 189 days (due to noisy adaptive lasso penalties).
			\end{itemize}
		\end{column}
	\end{columns}
	\refs{\cn{13} Shimizu et al.\ (2011); \cn{14} Granger (1969).\quad Data: CRSP via WRDS, FRED, Yahoo Finance}
\end{frame}

%----------------------------------------------------------------------------------------
%	4. METHODOLOGY II — INCORPORATING EDGE DIRECTIONS
%----------------------------------------------------------------------------------------

\begin{frame}
	\frametitle{Incorporating asymmetric edge orientations into hierarchical allocators}
	\dense
	\begin{columns}[T]
		\begin{column}{0.50\linewidth}
			\textbf{Structural equation model}\cn{15}
			\begin{center}\footnotesize
				$(I - M^{\top})\,x = \varepsilon$\\[4pt]
				$T \equiv (I - M^{\top})^{-1} = \sum_{k \ge 0} (M^{\top})^{k}$ \quad (exact, as $M$ nilpotent)\\[4pt]
				$\Sigma_{\mathrm{SEM}} = D_\sigma \,\big(T\,\Sigma_\varepsilon\,T^{\top}\big)\, D_\sigma,$\\[4pt]
                $\Sigma_\varepsilon = \mathrm{diag}(\sigma^2_{\varepsilon,i}),\ D_\sigma = \mathrm{diag}(\sigma_i)$
			\end{center}
            \vspace{-1em}
			\begin{itemize}
				\setlength{\itemsep}{3pt}
				\item $T_{ij}$ of a unit shock to asset $j$ on asset $i$ along every path\cn{16}; edge orientation information enters through its asymmetry.
				\item $\Sigma_\varepsilon$ is diagonal by design to force the model to explain dependence between assets via $T$, not sample correlation.
				\item Independent shocks remove the common market factor; the mechanism controls isolate this effect.
			\end{itemize}
		\end{column}
		\begin{column}{0.5\linewidth}
			\textbf{Two entry points, four allocators}\\[1em]
			{\scriptsize\adjustbox{max width=\linewidth}{
			\begin{tabular}{@{}l l l@{}}
				\toprule
				\textbf{ordering from} & \textbf{bisects on sample} & \textbf{bisects on $\Sigma_{\mathrm{SEM}}$} \\
				\midrule
				skeleton dendrogram & \textcolor{skel}{\skl{}} & \textcolor{orient}{\sem} \\
				topological sort & \textcolor{orient}{\topo} & \textcolor{orient}{\toposem} \\
				\bottomrule
			\end{tabular}}}\\[1em]
			\begin{itemize}
				\setlength{\itemsep}{3pt}
				\item Allocators are long-only, deterministic, same $M$.
			\end{itemize}
		\end{column}
	\end{columns}
	\refs{\cn{15} Bollen (1989); \cn{16} Wright (1934)}
\end{frame}

%----------------------------------------------------------------------------------------
%	5. METHODOLOGY III — DECOMPOSING THE EFFECT
%----------------------------------------------------------------------------------------

\begin{frame}
	\frametitle{Decomposing the effect: ablation design and controls to identify mechanisms}
	\dense
	\begin{columns}[T]
		\begin{column}{0.60\linewidth}
			\begin{tikzpicture}[
				rung/.style={draw=ICLBlue, rounded corners=2pt, text width=3.0cm, align=center, minimum height=0.6cm, font=\footnotesize, fill=white, inner sep=3pt},
				ctl/.style={draw=ctrl, dashed, rounded corners=2pt, text width=4.5cm, align=left, font=\scriptsize, fill=ICLLightGrey!50, inner sep=3pt},
				arr/.style={-{Stealth[length=4pt]}, thick},
				lbl/.style={font=\footnotesize\bfseries, fill=white, inner sep=1pt}]
                \node[rung] (r1) at (0,4.2) {\corr\\ \scriptsize graph-blind baseline};
                \node[rung] (r2) at (0,2.1) {\skl\\ \scriptsize skeleton-only};
                \node[rung] (r3) at (0,0)   {\sem{} / \topo{} / \toposem\\ \scriptsize orientation-aware};
                \draw[arr, skel]   (r1) -- node[lbl, text=skel]   {skeleton effect}    (r2);
                \draw[arr, orient] (r2) -- node[lbl, text=orient] {orientation effect} (r3);
                \node[ctl, anchor=west] (c3) at (2.15,3.4) {\pc{}: partial-correlation network with the DAG's edge count, no discovery $\to$ would any sparse network do?};
                \node[ctl, anchor=west] (c1) at (2.15,1.6) {\lw{}: same dendrogram, Ledoit--Wolf covariance\cn{17} $\to$ is it shrinkage?};
                \node[ctl, anchor=west] (c2) at (2.15,0.2) {\dfc{}: same dendrogram, covariance of residuals on the equal-weight market\cn{18} $\to$ is it market-factor removal?};
                \draw[ctrl] (r2.east) -- (1.9,2.1) -- (1.9,3.4) -- (c3.west);
                \draw[ctrl] (r3.east) -- (1.9,0) -- (1.9,1.6) -- (c1.west);
                \draw[ctrl] (1.9,0.2) -- (c2.west);
			\end{tikzpicture}\\[4pt]
			{\scriptsize Total DAG effect = skeleton + orientation.\\
            Robustness checks: \und{}; VARLiNGAM and ridge-Granger graphs; HERC; sparsity and cost sweeps.}
		\end{column}
		\begin{column}{0.38\linewidth}
			\textbf{Measurement}
			\begin{itemize}
				\setlength{\itemsep}{3pt}
				\item Pairwise comparisons with net Sharpe, 5 bps one-way transaction costs; stationary block bootstrap\cn{19}.
				\item Deflated Sharpe\cn{9,20} over all 146 configurations to account for multiplicity.
				\item Hansen's SPA\cn{10,21} and MCS\cn{11} for testing across whole groups (``families'')
				\item Power: minimum detectable effect $\approx 0.047$ (total), $0.037$ (orientation).
			\end{itemize}
		\end{column}
	\end{columns}
	\refs{\cn{17} Ledoit \& Wolf (2004); \cn{18} Sharpe (1963); \cn{19} Politis \& Romano (1994); \cn{20} Bailey \& López de Prado (2012); \cn{21} White (2000)}
\end{frame}

%----------------------------------------------------------------------------------------
%	6. IMPLEMENTATION — THE CODEBASE
%----------------------------------------------------------------------------------------
% Timing 10:40-12:00. Walk the diagram left to right: it is the data flow, and each stage is a package. Points to
% land: every allocator reads the same cached graph (the ablation depends on it); the whole thing is deterministic;
% vendored libraries are untouched; 88 tests and one reproduction script. Then the bug in two sentences: found by
% inspecting the baseline's dendrogram; fixing it turned an apparent gain into the null.

\begin{frame}
	\frametitle{Implementation}
	\dense
	\begin{center}
	\begin{tikzpicture}[
    	stg/.style={draw=ICLBlue, rounded corners=2pt, fill=ICLLightGrey!50, text width=2.8cm, align=left, font=\scriptsize, inner sep=3pt, minimum height=1.7cm, anchor=north west},
    	rail/.style={draw=ICLBlue, rounded corners=2pt, fill=ICLCream, align=left, font=\scriptsize, inner sep=4pt, text width=13.7cm, anchor=north west},
    	loop/.style={draw=ctrl, dashed, rounded corners=3pt, inner sep=4pt},
    	arr/.style={-{Stealth[length=4pt]}, ICLBlue, thick},
    	lab/.style={font=\tiny, text=ICLBlue, fill=white, inner sep=1pt, align=center}]
    	% --- 1: the panel, built once ----------------------------------------------------
    	\node[rail] (data) at (0,0) {132-column panel (\texttt{pipeline/data/}): 99 equities (CRSP prices via WRDS) and 33 drivers (FRED, Yahoo)};
    	% --- 2-5: the loop that runs at every rebalance ------------------------------------
    	\node[stg] (s2) at ($(data.south west)+(0.22,-0.5)$) {Extract the matrix's rows corresponding to the estimation window, and standardise.\par\texttt{closed\_loop.py}};
    	\node[stg] (s3) at ($(s2.north east)+(0.75,0)$) {Causal discovery, forbidding asset$\to$driver edges; keep the $99{\times}99$ asset block $M$.\par \texttt{pipeline/discovery/}};
    	\node[stg] (s4) at ($(s3.north east)+(0.75,0)$) {Allocation: \texttt{hrp}/\texttt{herc} cluster and bisect; \texttt{directed} builds $\Sigma_{\mathrm{SEM}}$ and controls; \texttt{topological} orders  DAG.\par \texttt{pipeline/portfolio/}};
    	\node[stg] (s5) at ($(s4.north east)+(0.75,0)$) {Record gross and net NAV over 21 trading days; 5 bps on one-way turnover.\par \texttt{portfolio/backtest.py}};
    	\draw[arr] (s2.east) -- node[lab] {window} (s3.west);
    	\draw[arr] (s3.east) -- node[lab] {graph} (s4.west);
    	\draw[arr] (s4.east) -- node[lab] {weights} (s5.west);
    	\draw[arr] (data.south -| s2) -- (s2.north);
    	\node[loop, fit=(s2)(s3)(s4)(s5)] (loopbox) {};
    	\node[font=\scriptsize, text=ctrl, anchor=south east, inner sep=1pt] at (loopbox.north east) {repeat every 21 trading days, Jan 2007 -- Nov 2024};
    	% --- 6-7: cells and collation ----------------------------------------------------
    	\node[rail] (cell) at ($(loopbox.south west)+(-0.22,-0.25)$) {Cells (discovery method, allocator, window, cost) saved in \texttt{results/<tag>/closed\_loop.pkl}; launched by \texttt{scripts/run\_phase\_ii.py}. 227 cells total, including sweeps, Phase I, and out-of-sample testing.\par
        Performance metrics, comparisons, and figures in \texttt{scripts/}.\par
        \texttt{reproduction.py} runs the complete pipeline. \texttt{causalnex}, \texttt{lingam}, and \texttt{nts-notears} are vendored.};
    \end{tikzpicture}
	\end{center}
	\vspace{-1em}
    {\scriptsize\textbf{\textcolor{ICLBlue}{Last-minute bugfix:}} All allocators accidentally had a nearest-PSD projection of its distance matrix, collapsing it to rank one. Removing it raised \corr{} by \pos{0.020} to \pos{0.046} and moved the graph-fed allocators by at most 0.005, creating an apparent performance gain.}
	\refs{Code and collated results: github.com/j-muthu/portfolio-optimisation-causal-networks. Exact reproduction needs WRDS/CRSP access; Yahoo prices are a fallback.}
\end{frame}

%----------------------------------------------------------------------------------------
%	7. RESULTS — NO GAIN, AND THE DECOMPOSITION
%----------------------------------------------------------------------------------------

\begin{frame}
	\frametitle{Results I: no overall performance gain, skeleton negative, orientation positive}
	\vspace{-4pt}
	\begin{columns}[T]
		\begin{column}{0.7\linewidth}
			% Left panel only: the right panel's point estimates are in the table below. Tune the trim fraction if the
			% panel boundary moves; alternatively export the left panel as figures/gradient_left.png and drop the trim.
			{\centering\adjincludegraphics[height=4.5cm, trim={0 0 {0.5\width} 0}, clip]{figures/phase_ii_window_gradient.png}\par}
            \vspace{-0.5em}
			{\scriptsize\adjustbox{max width=\linewidth}{\setlength{\tabcolsep}{2.5pt}
			\begin{tabular}{@{}l r r r r@{}}
				\toprule
				$\Delta$Sharpe \pv{$p$}\cn{19} & 189d & 252d & 378d & 504d \\
				\midrule
				\textcolor{skel}{\textbf{skeleton}} & \mn{0.019} \pv{0.21} & \mn{0.028} \pv{0.056} & \mn{0.003} \pv{0.80} & \mn{0.008} \pv{0.56} \\
				\textcolor{orient}{\textbf{orientation}}, alloc.\ step & \pos{0.015} \pv{0.22} & \pos{0.011} \pv{0.35} & \pos{0.009} \pv{0.45} & \pos{0.021} \pv{0.051} \\
				\textcolor{orient}{\textbf{orientation}}, both steps & \pos{0.016} \pv{0.31} & \pos{0.007} \pv{0.61} & \pos{0.015} \pv{0.33} & \pos{0.024} \pv{0.086} \\
				\textbf{total} & \mn{0.004} \pv{0.82} & \mn{0.017} \pv{0.36} & \pos{0.006} \pv{0.73} & \pos{0.013} \pv{0.44} \\
				\bottomrule
			\end{tabular}}}
		\end{column}
		\begin{column}{0.3\linewidth}
			\dense
			\begin{itemize}
				\setlength{\itemsep}{5pt}
				\item Only 6/20 DYNOTEARS cells beat HRP; 1/20 for VARLiNGAM.
				\item All 12 orientation vs. skeleton contrasts are non-negative but insignificant.
                \item Under VARLiNGAM graphs, the orientation contrasts show no consistent sign.
				\item Against the stricter skeleton allocator, the positive orientation effect falls.
			\end{itemize}
		\end{column}
	\end{columns}
	\refs{All contrasts on DYNOTEARS graphs.}
\end{frame}

%----------------------------------------------------------------------------------------
%	7. RESULTS — MECHANISM CONTROLS
%----------------------------------------------------------------------------------------
\begin{frame}
	\frametitle{Results II: neither effect needs a graph}
	\dense
{\scriptsize\adjustbox{max width=\linewidth}{
	\begin{tabular}{@{}l l r r r r c@{}}
		\toprule
		control & tests whether$\dots$ & 189 & 252 & 378 & 504 & reproduces it? \\
		\midrule
		\lw\cn{17}      & shrinkage explains the \textcolor{orient}{orientation} gain 
		               & \mn{0.002} & \mn{0.005} & \mn{0.022} & \mn{0.019} & \textbf{no} \\
		               & (\lw $-$ \skl) & \pval{0.86} & \pval{0.78} & \pval{0.19} & \pval{0.34} & \\[2pt]
		\dfc\cn{18,22} & market-factor removal explains the \textcolor{orient}{orientation} gain
		               & \pos{0.019} & \pos{0.010} & \pos{0.020} & \pos{0.027} & \textbf{yes} \\
		               & (\dfc $-$ \skl) & \pval{0.49} & \pval{0.69} & \pval{0.42} & \pval{0.26} & \\[2pt]
		\pc            & any sparse network reproduces the \textcolor{skel}{skeleton} cost
		               & \mn{0.020} & \pos{0.012} & \mn{0.005} & \pos{0.007} & \textbf{yes} \\
		               & (\pc $-$ \skl) & \pval{0.073} & \pval{0.23} & \pval{0.62} & \pval{0.48} & \\
		\bottomrule
	\end{tabular}}\\[2pt]
	\textcolor{ctrl}{$\Delta$Sharpe and two-sided block-bootstrap $p$-values.}}
    \vspace{-0.7em}
	\begin{columns}[T]
		\begin{column}{0.48\linewidth}
            \begin{itemize}
                \setlength{\itemsep}{3pt}
                \item Shrinkage does not explain orientation effect, as \lw{} underperforms \sem{}.
                \item Market-factor removal explains orientation effect, as \dfc{} matches \sem{} ($\Sigma_{\mathrm{SEM}}$'s independent shocks remove the common market factor).
                \item The DAG's density and sparsity alone explain the skeleton's effect, as \pc{} matches \skl{}.
            \end{itemize}
		\end{column}
		\begin{column}{0.48\linewidth}
			\begin{itemize}
				\setlength{\itemsep}{3pt}
				\item Ridge-Granger\cn{14} keeps directions but degrades magnitudes, which weakens the allocators bisecting on $\Sigma_{\mathrm{SEM}}$ -- direction alone does not drive performance.
				\item Market-factor removal driving orientation gain mirrors network-portfolio findings on centrality of assets\cn{23,24}.
			\end{itemize}
		\end{column}
	\end{columns}
	\refs{\cn{22} Laloux et al.\ (1999); \cn{23} Pozzi et al.\ (2013); \cn{24} Peralta \& Zareei (2016)}
\end{frame}

%----------------------------------------------------------------------------------------
%	7. RESULTS — INFERENCE, OUT OF SAMPLE, ROBUSTNESS
%----------------------------------------------------------------------------------------

\begin{frame}
	\frametitle{Results III: nothing survives the search; out of sample results are mixed}
	\dense
	\begin{columns}[T]
		\begin{column}{0.50\linewidth}
			\textbf{Hansen's SPA $p$-values}\cn{10}\\[3pt]
			{\scriptsize\adjustbox{max width=\linewidth}{\setlength{\tabcolsep}{3pt}
			\begin{tabular}{@{}l l r r r r r@{}}
				\toprule
				family & set & 189 & 252 & 378 & 504 & pooled \\
				\midrule
				graph-fed vs baseline & full (7) & 0.440 & 0.785 & 0.346 & 0.235 & 0.823 \\
				 & narrowed (5) & 0.673 & 0.886 & 0.273 & 0.213 & 0.827 \\
				oriented vs skeleton & full & 0.217 & 0.299 & 0.328 & 0.173 & 0.451 \\
				 & narrowed (3) & 0.361 & 0.333 & 0.175 & 0.076 & 0.399 \\
				\bottomrule
			\end{tabular}}}\\[6pt]
			\begin{itemize}
				\setlength{\itemsep}{3pt}
				\item Similarly, allocators' Deflated Sharpe ratios\cn{9} never reach 5\% significance.
                \item The 90\% Model Confidence Set\cn{11} keeps 16 of 17 strategies and drops only HSP.
				\item Point estimates are under 1.5$\times$ the bootstrap SEs; observed effects are about half the MDE.
			\end{itemize}
		\end{column}
		\begin{column}{0.46\linewidth}
			\textbf{Out of sample testing, Jan 2025 -- Aug 2026}\\[3pt]
			{\scriptsize\adjustbox{max width=\linewidth}{\setlength{\tabcolsep}{3pt}
			\begin{tabular}{@{}l r r r r@{}}
				\toprule
				$\Delta$Sharpe & 189 & 252 & 378 & 504 \\
				\midrule
				\textcolor{orient}{orientation} & \pos{0.035} & \pos{0.050} & \pos{0.105} & \pos{0.036} \\
				\textcolor{skel}{skeleton} & \pos{0.107} & \pos{0.071} & \mn{0.009} & \pos{0.100} \\
				\sem{} $-$ \dfc & \mn{0.090} & \mn{0.001} & \pos{0.005} & \mn{0.001} \\
				\bottomrule
			\end{tabular}}}\\[2pt]
            \begin{itemize}
				\setlength{\itemsep}{3pt}
                \item No statistically significant results.
                \item Orientation positive as in sample; the skeleton's contribution flips to positive at 3/4 windows.
            \end{itemize}

			\textbf{Robustness checks}
			\begin{itemize}
				\setlength{\itemsep}{3pt}
				\item Using HERC\cn{7} as a framework instead of HRP reproduces the negative skeleton and positive orientation effects; still insignificant.
				\item Howard et al.\cn{6} note causal DAGs may help in stressed regimes, but I find graph-fed allocators underperform HRP/HERC in recessions.
			\end{itemize}
		\end{column}
	\end{columns}
	\refs{}
\end{frame}

%----------------------------------------------------------------------------------------
%	8. CONCLUSIONS
%----------------------------------------------------------------------------------------

\begin{frame}
	\frametitle{Conclusions: replacing correlations with DAGs does not improve performance}
	\dense
	\begin{columns}[T]
		\begin{column}{0.32\linewidth}
			\textbf{Contributions}
			\begin{enumerate}
				\setlength{\itemsep}{3pt}
              \item An ablation design with graph-free mechanism
                    controls, fed by a reproducible DAG.
              \item Asymmetric inputs into hierarchical allocation:
                    $\Sigma_{\mathrm{SEM}}$ at the allocation step,
                    topological order at the ordering step.
              \item The finding: no gain at any window; the two components
                    are insignificant, offset each other, and are
                    reproducible without a graph.
			\end{enumerate}
		\end{column}
		\begin{column}{0.32\linewidth}
			\textbf{Implications}
			\begin{itemize}
				\setlength{\itemsep}{3pt}
				\item Market-factor removal performs as well as causal discovery (i.e. as well as standard HRP).
				\item Causal allocators should be tested via ablation and mechanism controls to identify the true performance effect drivers.
				\item Claims about ``causality'' improving prediction and allocation\cn{25} should be treated with skepticism!
			\end{itemize}
		\end{column}
		\begin{column}{0.32\linewidth}
			\textbf{Limitations}
			\begin{itemize}
				\setlength{\itemsep}{3pt}
				\item Underpowered: effects $\le 0.03$ vs MDE $0.04$--$0.05$.
				\item Orientation effect is DYNOTEARS-specific; long-only US large caps only.
				\item Three disclosed post hoc changes, included in the search-space correction.
				\item $R^2$-sortability\cn{26} and $\lambda$ untested.
			\end{itemize}
			\textbf{Future work:} NTS-NOTEARS\cn{27}; risk parity over structural shocks.
		\end{column}
	\end{columns}
	\vspace{-0.25em}
	\setbeamercolor{block body}{bg=ICLBlue, fg=white}
	\begin{block}{}
		\ImperialSansSemiBold The graph's skeleton is a worse clustering input than correlation, and its orientations only help through a mechanism that needs no graph.
	\end{block}
	\refs{\cn{25} López de Prado (2023); \cn{26} Reisach et al.\ (2023); \cn{27} Sun et al.\ (2023)}
\end{frame}

%----------------------------------------------------------------------------------------
%	CLOSING
%----------------------------------------------------------------------------------------

\begingroup
	\setbeamercolor{background canvas}{bg=ICLBlue} % Slide background color
	\setbeamercolor{normal text}{fg=white}\usebeamercolor[fg]{normal text} % Slide text color
	\setbeamertemplate{closing slide logo}[logo]{ICL_Logo_White.pdf} % Imperial logo color, use 'ICL_Logo_White.pdf' for white and 'ICL_Logo_Blue.pdf' for blue
	\setbeamertemplate{closing slide text}[text]{Thank you} % Slide text
	\usebeamertemplate{closing slide} % Output the closing slide
\endgroup

\end{document}